% Sezione 2: Robot Mobili e Algoritmi di Navigazione
\section{Stato dell'arte: robot mobili, navigazione e protocolli}

% Slide 1: Classificazione
\begin{frame}
    \frametitle{Classificazione dei Robot Mobili}

    \begin{block}{AMR vs AGV}
        \begin{itemize}
            \item \textbf{AGV}: Percorsi fissi, infrastruttura dedicata, bassa flessibilità.
            \item \textbf{AMR}: Navigazione autonoma, adattabilità, nessuna infrastruttura.
        \end{itemize}
    \end{block}


    \textbf{Tipologie di Locomozione}
    \begin{itemize}
        \item \textbf{Ruote}: Differential drive, omnidirezionali.
        \item \textbf{Cingoli}: Terreni irregolari.
        \item \textbf{Zampe}: Ambienti complessi.
        \item \textbf{Volanti}: Droni/UAV.
    \end{itemize}

\end{frame}

% Slide 2: Algoritmi di Navigazione
\begin{frame}
    \frametitle{Algoritmi di Navigazione}
    \begin{columns}
        \column{0.5\textwidth}
        \textbf{Path Planning (Globale)}
        \begin{itemize}
            \item \textbf{A*}: Ricerca ottima su griglia.
            \item \textbf{Dijkstra}: Ricerca esaustiva.
            \item \textbf{RRT}: Esplorazione probabilistica.
        \end{itemize}

        \column{0.5\textwidth}
        \textbf{Obstacle Avoidance (Locale)}
        \begin{itemize}
            \item \textbf{VFH}: Istogramma polare ostacoli.
            \item \textbf{DWA}: Ottimizzazione velocità ($v, \omega$).
            \item \textbf{Potential Fields}: Campi di forze.
        \end{itemize}
    \end{columns}
    \begin{block}{Localizzazione}
        Tecniche per stimare la posizione del robot nell'ambiente.
        \begin{itemize}
            \item \textbf{Odometry}: Basata su encoder e IMU, soggetta a deriva.
            \item \textbf{Beaconing}: Utilizzo di segnali esterni (es. RFID).
            \item \textbf{SLAM}: Simultaneous Localization and Mapping.
        \end{itemize}
    \end{block}
\end{frame}

% Slide 3: SLAM e FAST-LIO 2
\begin{frame}
    \frametitle{SLAM: FAST-LIO 2}
    \textbf{Problema}: Stimare posa e mappa simultaneamente in assenza di GPS.

    \vspace{0.5cm}
    \begin{block}{Perché FAST-LIO 2?}
        Algoritmo LiDAR-Inertial Odometry (LIO) allo stato dell'arte.
        \begin{itemize}
            \item \textbf{Robustezza}: Funziona in ambienti spogli (corridoi lunghi) grazie alla \textit{registrazione diretta} (no feature extraction).
            \item \textbf{Efficienza}: Struttura \textbf{ikd-Tree} per aggiornamento mappa veloce ($>100$ Hz).
            \item \textbf{Accuratezza}: Fusione sensoriale (tightly-coupled) LiDAR + IMU.
            \item \textbf{Testato}: Già validato sul sensore acquistato (Livox Mid-360).
        \end{itemize}
    \end{block}
\end{frame}

% Slide 4: Architettura FAST-LIO 2
\begin{frame}
    \frametitle{Architettura FAST-LIO 2}
    \begin{figure}[htbp]
        \centering
        \includegraphics[width=\textwidth]{2/fastlio2_block_diagram.png}
        \caption{Schema semplificato: IMU $\to$ Predizione $\to$ Correzione (LiDAR) $\to$ Mappa.}
    \end{figure}
\end{frame}

% Slide 5: Protocolli Industriali - bus CAN
\begin{frame}
    \frametitle{Protocolli di Comunicazione: CAN Bus}
    \textbf{Controller Area Network (CAN)}: bus seriale multi-master per autoveicoli.
    \begin{block}{Caratteristiche Principali}
        \begin{itemize}
            \item \textbf{Architettura Multi-Master}: Tutti i nodi possono trasmettere.
            \item \textbf{Arbitraggio Non-Destruttivo}: Priorità basata su ID messaggio.
            \item \textbf{Rilevamento Errori}: CRC, bit stuffing, monitoraggio ACK.
            \item \textbf{Velocità}: Fino a 1 Mbps su brevi distanze.
        \end{itemize}
    \end{block}

    \begin{block}{Applicazioni Industriali}
        Utilizzato in automotive, automazione industriale, robotica mobile.
    \end{block}
\end{frame}
% Slide 5: Protocolli Industriali - CANopen
\begin{frame}
    \frametitle{Protocolli di Comunicazione: CANopen}
    Protocollo applicativo basato su bus CAN per atomotive e automazione industriale.
    \begin{block}{Struttura Dati}
        \begin{itemize}
            \item \textbf{OD (Object Dictionary)}: Mappa parametri (Index/Subindex).
            \item \textbf{PDO}: Dati real-time (es. velocità, encoder).
            \item \textbf{SDO}: Configurazione (confermato).
        \end{itemize}
    \end{block}

    \begin{block}{Gestione Rete}
        \begin{itemize}
            \item \textbf{NMT}: Stati (Operational, Stopped).
            \item \textbf{Heartbeat}: Monitoraggio nodi.
            \item \textbf{Emergency}: Segnalazione guasti.
        \end{itemize}
    \end{block}
\end{frame}

% Slide 6: Profilo CiA 402
\begin{frame}
    \frametitle{Estensioni CANopen: CiA 402}
    
    Estensione del protocollo per il controllo di azionamenti elettrici.

    \begin{itemize}
        \item \textbf{Macchina a Stati}: Gestione sicurezza e abilitazione potenza.
        \item \textbf{Modalità Operative}:
              \begin{itemize}
                  \item \textit{Profile Velocity}: Controllo velocità.
                  \item \textit{Profile Position}: Controllo posizione.
                  \item \textit{Profile Torque}: Controllo coppia.
              \end{itemize}
    \end{itemize}

    \begin{alertblock}{Vantaggi}
        Determinismo, affidabilità e interoperabilità tra produttori diversi.
    \end{alertblock}
\end{frame}
